元类（metaclass）常常出现在列举 Python 独特特性的清单中，但鲜有人理解它们在实践中究竟完成了什么。metaclass 这个名字隐约暗示了一种超越类的概念。简而言之，元类让你能够拦截 Python 的 \texttt{class} 语句，并在每次定义一个类时提供特殊行为。

同样神秘的是 Python 用于动态定制属性（attribute）访问的内置特性。配合 Python 的面向对象构造，这些设施提供了强大的工具，让从简单的类过渡到复杂的类变得更加平滑。

然而，这些能力也伴随着诸多陷阱。动态属性让你能够覆盖对象并引发意料之外的副作用。元类能制造出令新手难以接近的极其怪异的行为。重要的是，你要遵循最小惊奇原则，并且只使用这些机制来实现那些已被充分理解的惯用法（idiom）。

\section{Item 58: 使用普通属性而非 setter 和 getter 方法}

从其他语言转到 Python 的程序员，可能会很自然地试图在自己的类中实现显式的 getter 和 setter 方法，用以访问受保护（protected）属性（相关背景见 Item 55「优先使用公有属性而非私有属性」）：

\begin{codebox}[python]
class OldResistor:
    def __init__(self, ohms):
        self._ohms = ohms
    def get_ohms(self):
        return self._ohms
    def set_ohms(self, ohms):
        self._ohms = ohms
\end{codebox}

使用这些 setter 和 getter 方法很简单，但这并不 Pythonic（地道）：

\begin{codebox}[python]
r0 = OldResistor(50e3)
print("Before:", r0.get_ohms())
r0.set_ohms(10e3)
print("After: ", r0.get_ohms())
>>>
Before: 50000.0
After:  10000.0
\end{codebox}

这类方法在像原地自增这样的操作时尤其笨拙：

\begin{codebox}[python]
r0.set_ohms(r0.get_ohms() - 4e3)
assert r0.get_ohms() == 6e3
\end{codebox}

不过，这些实用方法确实有助于定义类的接口，让封装功能、校验用法、划定边界变得更容易。在设计类以确保随着时间推移而演进时不会破坏调用方这一点上，这些都是重要目标。

然而在 Python 中，你永远不需要实现像这样的显式 setter 或 getter 方法。相反，你应该始终以简单的公有属性来开始你的实现，就像我这里做的：

\begin{codebox}[python]
class Resistor:
    def __init__(self, ohms):
        self.ohms = ohms
        self.voltage = 0
        self.current = 0
r1 = Resistor(50e3)
r1.ohms = 10e3
\end{codebox}

这些属性让像原地自增这样的操作变得自然又清晰：

\begin{codebox}[python]
r1.ohms += 5e3
\end{codebox}

后来，如果我决定需要在设置某个属性时附加特殊行为，我可以迁移到 \texttt{@property} 装饰器（相关背景见 Item 38「用 \texttt{functools.wraps} 定义函数装饰器」）及其对应的 \texttt{setter} 属性。这里，我定义了一个 \texttt{Resistor} 的子类，让我可以通过赋值 \texttt{voltage} 属性来改变 \texttt{current}。注意，为了让这段代码正确工作，setter 和 getter 方法的名字必须与预期的属性名（\texttt{voltage}）一致：

\begin{codebox}[python]
class VoltageResistance(Resistor):
    def __init__(self, ohms):
        super().__init__(ohms)
        self._voltage = 0
    @property
    def voltage(self):
        return self._voltage
    @voltage.setter
    def voltage(self, voltage):
        self._voltage = voltage
        self.current = self._voltage / self.ohms
\end{codebox}

现在，赋值 \texttt{voltage} 属性会运行 \texttt{voltage} 的 setter 方法，该方法进而会更新对象的 \texttt{current} 属性以保持同步：

\begin{codebox}[python]
r2 = VoltageResistance(1e2)
print(f"Before: {r2.current:.2f} amps")
r2.voltage = 10
print(f"After:  {r2.current:.2f} amps")
>>>
Before: 0.00 amps
After:  0.10 amps
\end{codebox}

为属性指定 setter 还能让我对传入类的值进行类型检查和校验。这里我定义一个类，确保所有的电阻值都高于零欧姆：

\begin{codebox}[python]
class BoundedResistance(Resistor):
    def __init__(self, ohms):
        super().__init__(ohms)
    @property
    def ohms(self):
        return self._ohms
    @ohms.setter
    def ohms(self, ohms):
        if ohms <= 0:
            raise ValueError(f"ohms must be > 0; got {ohms}")
        self._ohms = ohms
\end{codebox}

现在，给该属性赋一个无效的电阻值会引发异常：

\begin{codebox}[python]
r3 = BoundedResistance(1e3)
r3.ohms = 0
>>>
Traceback ...
ValueError: ohms must be > 0; got 0
\end{codebox}

如果我向构造函数传入一个无效值，同样会引发异常：

\begin{codebox}[python]
BoundedResistance(-5)
>>>
Traceback ...
ValueError: ohms must be > 0; got -5
\end{codebox}

这是因为 \texttt{BoundedResistance.\_\_init\_\_} 调用了 \texttt{Resistor.\_\_init\_\_}，后者赋值 \texttt{self.ohms = -5}。该赋值使得 \texttt{BoundedResistance} 的 \texttt{@ohms.setter} 方法被调用，而它在对象构造完成之前就立即运行了校验代码。

我甚至可以使用 \texttt{@property} 把父类的属性变成不可变的（另一种做法见 Item 56「优先使用 \texttt{dataclasses} 创建不可变对象」）：

\begin{codebox}[python]
class FixedResistance(Resistor):
    def __init__(self, ohms):
        super().__init__(ohms)
    @property
    def ohms(self):
        return self._ohms
    @ohms.setter
    def ohms(self, ohms):
        if hasattr(self, "_ohms"):
            raise AttributeError("Ohms is immutable")
        self._ohms = ohms
\end{codebox}

在构造之后尝试给该属性赋值会引发异常：

\begin{codebox}[python]
r4 = FixedResistance(1e3)
r4.ohms = 2e3
>>>
Traceback ...
AttributeError: Ohms is immutable
\end{codebox}

当你使用 \texttt{@property} 方法来实现 setter 和 getter 时，要确保你实现的行为不会令人意外。例如，不要在 getter 属性方法中设置其他属性：

\begin{codebox}[python]
class MysteriousResistor(Resistor):
    @property
    def ohms(self):
        self.voltage = self._ohms * self.current
        return self._ohms
    @ohms.setter
    def ohms(self, ohms):
        self._ohms = ohms
\end{codebox}

在 getter 属性方法中设置其他属性会导致极其怪异的行为：

\begin{codebox}[python]
r7 = MysteriousResistor(10)
r7.current = 0.1
print(f"Before: {r7.voltage:.2f}")
r7.ohms
print(f"After:  {r7.voltage:.2f}")
>>>
Before: 0.00
After:  1.00
\end{codebox}

最好的做法是只在 \texttt{@property.setter} 方法中修改相关的对象状态。同时务必避免调用方可能没预料到的、超出对象之外的任何其他副作用，例如动态导入模块、运行缓慢的辅助函数、执行 I/O，或者进行昂贵的数据库查询。类的使用者会期望其属性像任何其他 Python 对象一样：快速而简便。任何更复杂或更慢的事情，都用普通方法来完成。

\texttt{@property} 最大的缺点是，某个属性的方法只能被子类共享。不相关的类无法共享同一份实现。不过，Python 还支持描述符（descriptor）（见 Item 60「用描述符实现可复用的 \texttt{@property} 方法」），它支持可复用的属性逻辑以及许多其他用例。

\subsection{需要记住的要点}

\begin{itemize}
\item 用简单的公有属性来定义新的类接口，避免定义 setter 和 getter 方法。
\item 当对象的属性被访问时，使用 \texttt{@property} 来定义特殊行为。
\item 遵循最小意外原则，避免在 \texttt{@property} 方法中出现奇怪的副作用。
\item 确保 \texttt{@property} 方法运行迅速；对于缓慢或复杂的工作——尤其是涉及 I/O 或引发副作用的工作——改用普通方法。
\end{itemize}

\section{Item 59: 考虑用 @property 而非重构属性}

内置的 \texttt{@property} 装饰器让对实例属性的简单访问能够表现得更智能（见 Item 58「使用普通属性而非 setter 和 getter 方法」）。\texttt{@property} 一个高级但常见的用途，是把曾经的一个简单数值属性转换成一个即时计算。这极其有用，因为它让你能够把某个类的所有现有用法迁移到新行为上，而无需重写任何调用点（如果你无法控制调用代码，这一点尤其重要）。\texttt{@property} 还为逐步改进接口提供了一个重要的过渡方案。

例如，假设我想用普通的 Python 对象实现一个漏桶（leaky-bucket）限速配额系统。这里，\texttt{Bucket} 类表示剩余的配额数量，以及配额可用的时长：

\begin{codebox}[python]
from datetime import datetime, timedelta
class Bucket:
    def __init__(self, period):
        self.period_delta = timedelta(seconds=period)
        self.reset_time = datetime.now()
        self.quota = 0
    def __repr__(self):
        return f"Bucket(quota={self.quota})"
\end{codebox}

漏桶算法的工作原理是：每当桶被填充时，确保配额数量不会从一个周期延续到下一个周期：

\begin{codebox}[python]
def fill(bucket, amount):
    now = datetime.now()
    if (now - bucket.reset_time) > bucket.period_delta:
        bucket.quota = 0
        bucket.reset_time = now
    bucket.quota += amount
\end{codebox}

每次配额消费者想要做某件事时，它都必须先确保能够扣减它需要使用掉的配额数量：

\begin{codebox}[python]
def deduct(bucket, amount):
    now = datetime.now()
    if (now - bucket.reset_time) > bucket.period_delta:
        return False  # Bucket hasn't been filled this period
    if bucket.quota - amount < 0:
        return False  # Bucket was filled, but not enough
    bucket.quota -= amount
    return True       # Bucket had enough, quota consumed
\end{codebox}

要使用这个类，首先我把桶填满：

\begin{codebox}[python]
bucket = Bucket(60)
fill(bucket, 100)
print(bucket)
>>>
Bucket(quota=100)
\end{codebox}

然后，我扣减我需要的配额：

\begin{codebox}[python]
if deduct(bucket, 99):
    print("Had 99 quota")
else:
    print("Not enough for 99 quota")
print(bucket)
>>>
Had 99 quota
Bucket(quota=1)
\end{codebox}

最终，我会被阻止继续推进，因为我尝试扣减的配额超过了可用量。在这种情况下，桶的配额水平保持不变：

\begin{codebox}[python]
if deduct(bucket, 3):
    print("Had 3 quota")
else:
    print("Not enough for 3 quota")
print(bucket)
>>>
Not enough for 3 quota
Bucket(quota=1)
\end{codebox}

这个实现的问题在于，我永远不知道桶最初有多少配额。配额在周期内被扣减，直到降为零。到那时，在桶被重新填充之前，\texttt{deduct} 将始终返回 \texttt{False}。当这种情况发生时，如果能知道 \texttt{deduct} 的调用方被挡住，是因为 \texttt{Bucket} 耗尽了配额，还是因为 \texttt{Bucket} 在本周期内一开始就没有配额，那会很有用。

为了修复这个问题，我可以修改这个类，使其记录本周期内发放的 \texttt{max\_quota} 以及同一周期内已消耗的 \texttt{quota\_consumed}：

\begin{codebox}[python]
class NewBucket:
    def __init__(self, period):
        self.period_delta = timedelta(seconds=period)
        self.reset_time = datetime.now()
        self.max_quota = 0
        self.quota_consumed = 0
    def __repr__(self):
        return (
            f"NewBucket(max_quota={self.max_quota}, "
            f"quota_consumed={self.quota_consumed})"
        )
\end{codebox}

为了与原 \texttt{Bucket} 类的既有接口保持一致，我用一个 \texttt{@property} 方法，借助这些新属性即时算出当前的配额水平：

\begin{codebox}[python]
    @property
    def quota(self):
        return self.max_quota - self.quota_consumed
\end{codebox}

当 \texttt{quota} 属性被赋值时，为了与 \texttt{fill} 和 \texttt{deduct} 函数对类的当前用法保持兼容，我会采取特殊动作：

\begin{codebox}[python]
    @quota.setter
    def quota(self, amount):
        delta = self.max_quota - amount
        if amount == 0:
            # Quota being reset for a new period
            self.quota_consumed = 0
            self.max_quota = 0
        elif delta < 0:
            # Quota being filled during the period
            self.max_quota = amount + self.quota_consumed
        else:
            # Quota being consumed during the period
            self.quota_consumed = delta
\end{codebox}

重新运行上面的示例代码会得出相同的结果：

\begin{codebox}[python]
bucket = NewBucket(60)
print("Initial", bucket)
fill(bucket, 100)
print("Filled", bucket)
if deduct(bucket, 99):
    print("Had 99 quota")
else:
    print("Not enough for 99 quota")
print("Now", bucket)
if deduct(bucket, 3):
    print("Had 3 quota")
else:
    print("Not enough for 3 quota")
print("Still", bucket)
>>>
Initial NewBucket(max_quota=0, quota_consumed=0)
Filled NewBucket(max_quota=100, quota_consumed=0)
Had 99 quota
Now NewBucket(max_quota=100, quota_consumed=99)
Not enough for 3 quota
Still NewBucket(max_quota=100, quota_consumed=99)
\end{codebox}

最妙的地方在于，使用 \texttt{Bucket.quota} 的代码无需改动，也无需知道这个类已经发生了变化。\texttt{Bucket} 的新用法可以正确地直接访问 \texttt{max\_quota} 和 \texttt{quota\_consumed}。

我特别喜欢 \texttt{@property}，因为它让你能够随着时间推移，逐步迈向更好的数据模型。读到上面的 \texttt{Bucket} 示例，你或许会想，\texttt{fill} 和 \texttt{deduct} 从一开始就应该被实现成实例方法。尽管你大概是对的，但在实际中，有很多情况下，对象最初带有定义含糊的接口，或者只是充当笨拙的数据容器（相关示例见 Item 51「优先使用 \texttt{dataclasses} 定义轻量级类」）。当代码随时间增长、作用域扩大、多位作者贡献却无人考虑长期可维护性时，就会发生这种情况。

\texttt{@property} 是一个能帮你解决真实代码中会遇到的问题的工具。不要过度使用它。当你发现自己反复地扩展 \texttt{@property} 方法时，大概就到了应该重构你的类、而不是继续粉饰你糟糕设计的时候了。

\subsection{需要记住的要点}

\begin{itemize}
\item 使用 \texttt{@property} 为已有的实例属性赋予新的功能。
\item 通过使用 \texttt{@property}，逐步迈向更好的数据模型。
\item 当你发现自己过度使用 \texttt{@property} 时，考虑重构这个类以及所有的调用点。
\end{itemize}

\section{Item 60: 使用描述符实现可复用的 @property 方法}

内置（built-in）的 \texttt{@property}（参见 Item 58「用普通属性代替 setter 和 getter 方法」与 Item 59「考虑用 \texttt{@property} 代替重构属性」）最大的问题在于复用（reuse）。它所装饰的方法无法在同一个类的多个属性之间复用，也无法被不相干的类复用。

例如，假设我想写一个类，用来验证学生在一次家庭作业（homework）中得到的成绩是一个百分比：

\begin{codebox}[python]
class Homework:
    def __init__(self):
        self._grade = 0
    @property
    def grade(self):
        return self._grade
    @grade.setter
    def grade(self, value):
        if not (0 <= value <= 100):
            raise ValueError("Grade must be between 0 and 100")
        self._grade = value
\end{codebox}

使用 \texttt{@property} 让这个类很好用：

\begin{codebox}[python]
galileo = Homework()
galileo.grade = 95
\end{codebox}

再假设我还想给这个学生一门考试（exam）的成绩，而考试包含多个科目，每个科目有各自独立的成绩：

\begin{codebox}[python]
class Exam:
    def __init__(self):
        self._writing_grade = 0
        self._math_grade = 0
    @staticmethod
    def _check_grade(value):
        if not (0 <= value <= 100):
            raise ValueError("Grade must be between 0 and 100")
\end{codebox}

这样很快会变得繁琐。对于考试的每一科，我都需要新增一个 \texttt{@property} 以及相关的验证：

\begin{codebox}[python]
    @property
    def writing_grade(self):
        return self._writing_grade
    @writing_grade.setter
    def writing_grade(self, value):
        self._check_grade(value)
        self._writing_grade = value
    @property
    def math_grade(self):
        return self._math_grade
    @math_grade.setter
    def math_grade(self, value):
        self._check_grade(value)
        self._math_grade = value
\end{codebox}

此外，这种做法并不通用。如果我想在家庭作业和考试之外，把这套百分比验证逻辑复用到其他类里，就不得不一遍又一遍地重复编写 \texttt{@property} 的样板代码和 \texttt{\_check\_grade} 方法。

在 Python 中，更好的做法是使用描述符（descriptor）。描述符协议（descriptor protocol）定义了语言如何解释属性访问。描述符类可以提供 \texttt{\_\_get\_\_} 和 \texttt{\_\_set\_\_} 方法，让你无需任何样板代码就能复用成绩验证逻辑。就这一点而言，描述符也比 mix-in（混入类）更胜一筹（参见 Item 54「用 mix-in 类组合功能」），因为它让你能在同一个类的多个不同属性上复用同一套逻辑。

下面我定义一个名为 \texttt{Exam} 的新类，其类属性是 \texttt{Grade} 的实例。\texttt{Grade} 类实现了描述符协议：

\begin{codebox}[python]
class Grade:
    def __get__(self, instance, instance_type):
        ...
    def __set__(self, instance, value):
        ...
class Exam:
    # 类属性
    math_grade = Grade()
    writing_grade = Grade()
    science_grade = Grade()
\end{codebox}

在解释 \texttt{Grade} 类如何工作之前，有一点很重要：要理解当在 \texttt{Exam} 实例上访问这样的描述符属性时，Python 会做什么。当我赋值一个属性时：

\begin{codebox}[python]
exam = Exam()
exam.writing_grade = 40
\end{codebox}

它会被解释为：

\begin{codebox}[python]
Exam.__dict__["writing_grade"].__set__(exam, 40)
\end{codebox}

当我读取一个属性时：

\begin{codebox}[python]
exam.writing_grade
\end{codebox}

它会被解释为：

\begin{codebox}[python]
Exam.__dict__["writing_grade"].__get__(exam, Exam)
\end{codebox}

驱动这一行为的是 \texttt{object} 的 \texttt{\_\_getattribute\_\_} 方法（背景知识参见 Item 61「用 \texttt{\_\_getattr\_\_}、\texttt{\_\_getattribute\_\_} 和 \texttt{\_\_setattr\_\_} 实现惰性属性」）。简而言之，当 \texttt{Exam} 实例没有名为 \texttt{writing\_grade} 的属性时，Python 会回退到 \texttt{Exam} 类的属性。如果这个类属性是一个带有 \texttt{\_\_get\_\_} 和 \texttt{\_\_set\_\_} 方法的对象，Python 就假定你想遵循描述符协议。

了解了这一行为，以及我在 \texttt{Homework} 类中如何用 \texttt{@property} 做成绩验证，下面尝试实现一个合理的 \texttt{Grade} 描述符：

\begin{codebox}[python]
class Grade:
    def __init__(self):
        self._value = 0
    def __get__(self, instance, instance_type):
        return self._value
    def __set__(self, instance, value):
        if not (0 <= value <= 100):
            raise ValueError("Grade must be between 0 and 100")
        self._value = value
\end{codebox}

遗憾的是，这是错的，会导致异常行为。在单个 \texttt{Exam} 实例上访问多个属性时，行为符合预期：

\begin{codebox}[python]
class Exam:
    math_grade = Grade()
    writing_grade = Grade()
    science_grade = Grade()
first_exam = Exam()
first_exam.writing_grade = 82
first_exam.science_grade = 99
print("Writing", first_exam.writing_grade)
print("Science", first_exam.science_grade)
>>>
Writing 82
Science 99
\end{codebox}

但是，在多个 \texttt{Exam} 实例上访问这些属性时，会出现令人意外的行为：

\begin{codebox}[python]
second_exam = Exam()
second_exam.writing_grade = 75
print(f"Second {second_exam.writing_grade} is right")
print( f"First  {first_exam.writing_grade} is wrong; "
f"should be 82")
>>>
Second 75 is right
First  75 is wrong; should be 82
\end{codebox}

问题在于，对于类属性 \texttt{writing\_grade}，一个单独的 \texttt{Grade} 实例被所有 \texttt{Exam} 实例共享。这个 \texttt{Grade} 实例在程序生命周期中只会被构造一次，也就是在 \texttt{Exam} 类首次定义时，而不是每次创建 \texttt{Exam} 实例时。

为了解决这个问题，我需要让 \texttt{Grade} 类为每个独立的 \texttt{Exam} 实例分别记录其值。办法是把每个实例的状态保存在一个字典里：

\begin{codebox}[python]
class DictGrade:
    def __init__(self):
        self._values = {}
    def __get__(self, instance, instance_type):
        if instance is None:
            return self
        return self._values.get(instance, 0)
    def __set__(self, instance, value):
        if not (0 <= value <= 100):
            raise ValueError("Grade must be between 0 and 100")
        self._values[instance] = value
\end{codebox}

这个实现简单且运行良好，但仍有一个陷阱：它会泄漏内存。\texttt{\_values} 字典持有对每个曾被传入 \texttt{\_\_set\_\_} 的 \texttt{Exam} 实例的引用，贯穿整个程序的生命周期。这导致这些实例的引用计数（reference count）永远不会降到零，从而阻止了垃圾回收器（garbage collector）的清理（如何检测这类问题，参见 Item 115「用 \texttt{tracemalloc} 理解内存使用与泄漏」）。

相反，你应该依赖 Python 描述符的 \texttt{\_\_set\_name\_\_} 特殊方法（背景知识参见 Item 64「用 \texttt{\_\_set\_name\_\_} 标注类属性」）。这个方法会在类定义完成后，对每个描述符实例调用。关键在于，赋给描述符实例的那个类属性的名字是由 Python 提供的。这样你就能计算出一个字符串，用作每个对象属性的名字（这里是一个以 \texttt{"\_"} 开头的受保护字段）：

\begin{codebox}[python]
class NamedGrade:
    def __set_name__(self, owner, name):
        self.internal_name = "_" + name
\end{codebox}

我可以借助描述符提供的 \texttt{internal\_name}，在对象上调用 \texttt{setattr} 和 \texttt{getattr}，来存储和取出相应的属性数据：

\begin{codebox}[python]
    def __get__(self, instance, instance_type):
        if instance is None:
            return self
        return getattr(instance, self.internal_name)
    def __set__(self, instance, value):
        if not (0 <= value <= 100):
            raise ValueError("Grade must be between 0 and 100")
        setattr(instance, self.internal_name, value)
\end{codebox}

现在，用这个改进版的描述符定义一个新类，可以看到描述符的属性数据存放在对象的实例字典（\texttt{\_\_dict\_\_}）中：

\begin{codebox}[python]
class NamedExam:
    math_grade = NamedGrade()
    writing_grade = NamedGrade()
    science_grade = NamedGrade()
first_exam = NamedExam()
first_exam.math_grade = 78
first_exam.writing_grade = 89
first_exam.science_grade = 94
print(first_exam.__dict__)
>>>
{'_math_grade': 78, '_writing_grade': 89, '_science_grade': 94}
\end{codebox}

与前面的实现不同，这个版本不会泄漏内存，因为当 \texttt{NamedExam} 对象被垃圾回收时，它所有的属性数据（包括由描述符赋的值）也都会被释放。

\subsection{需要记住的要点}

\begin{itemize}
\item 通过定义自己的描述符类，复用 \texttt{@property} 方法的行为与验证。
\item 配合使用 \texttt{\_\_set\_name\_\_}、\texttt{setattr} 和 \texttt{getattr}，把描述符所需的数据存放在对象的实例字典中，以避免内存泄漏。
\item 不要纠结于去精确理解 \texttt{\_\_getattribute\_\_} 是如何借助描述符协议来读取和设置属性的。
\end{itemize}

\section{Item 61: 使用 \texttt{\_\_getattr\_\_}、\texttt{\_\_getattribute\_\_} 和 \texttt{\_\_setattr\_\_} 实现惰性属性}

Python 的 \texttt{object} 基类提供了一些钩子（hook），使得编写用于粘合各系统的通用代码变得很容易。例如，假设我想把数据库中的记录表示为 Python 对象。我假定数据库的模式（schema）已经在别处定义好了。在大多数语言中，我需要显式地在代码中说明数据库模式如何映射到程序里的类和对象。然而在 Python 中，我可以在运行时（run-time）以通用的方式完成这种对象关系映射（object-relational mapping），因此无需任何样板代码。

这怎么可能实现呢？普通的实例属性、\texttt{@property} 方法和描述符（descriptor）都无法做到这一点，因为它们都需要事先定义好。Python 通过 \texttt{\_\_getattr\_\_} 这个特殊方法（special method）实现了这种动态行为。如果某个类定义了 \texttt{\_\_getattr\_\_}，那么每当在对象的实例字典中找不到某个属性时，Python 就会调用该方法。这里我定义一个 \texttt{\_\_getattr\_\_} 钩子，它会把一个属性插入对象的实例字典，以此证明它确实运行过：

\begin{codebox}[python]
class LazyRecord:
    def __init__(self):
        self.exists = 5
    def __getattr__(self, name):
        value = f"Value for {name}"
        setattr(self, name, value)
        return value
\end{codebox}

例如，当我访问对象上不存在的属性 \texttt{foo} 时，Python 会调用上面的 \texttt{\_\_getattr\_\_} 方法，该方法会修改实例字典 \texttt{\_\_dict\_\_}：

\begin{codebox}[python]
data = LazyRecord()
print("Before:", data.__dict__)
print("foo:   ", data.foo)
print("After: ", data.__dict__)
>>>
Before: {'exists': 5}
foo:    Value for foo
After:  {'exists': 5, 'foo': 'Value for foo'}
\end{codebox}

我可以为 \texttt{LazyRecord} 添加日志（logging），以显示 \texttt{\_\_getattr\_\_} 实际被调用的时机。请注意，在这个实现中我调用了 \texttt{super().\_\_getattr\_\_()}，以便使用父类的 \texttt{\_\_getattr\_\_} 实现来获取真实的属性值，从而避免无限递归（详见 Item 53（"使用 \texttt{super} 初始化父类"））：

\begin{codebox}[python]
class LoggingLazyRecord(LazyRecord):
    def __getattr__(self, name):
        print(
            f"* Called __getattr__({name!r}), "
            f"populating instance dictionary"
        )
        result = super().__getattr__(name)
        print(f"* Returning {result!r}")
        return result
data = LoggingLazyRecord()
print("exists:     ", data.exists)
print("First foo:  ", data.foo)
print("Second foo: ", data.foo)
>>>
exists:      5
* Called __getattr__('foo'), populating instance dictionary
* Returning 'Value for foo'
First foo:   Value for foo
Second foo:  Value for foo
\end{codebox}

\texttt{exists} 属性存在于实例字典中，因此永远不会为它调用 \texttt{\_\_getattr\_\_}。\texttt{foo} 属性起初不在实例字典中，所以第一次会调用 \texttt{\_\_getattr\_\_}。但是，为 \texttt{foo} 调用 \texttt{\_\_getattr\_\_} 时，它还会执行一次 \texttt{setattr}，从而把 \texttt{foo} 填入实例字典。这就是为什么第二次访问 \texttt{foo} 时，不会记录对 \texttt{\_\_getattr\_\_} 的调用。

这种行为对于像惰性访问无模式（schemaless）数据这样的用例特别有用。\texttt{\_\_getattr\_\_} 会运行一次，完成加载属性的繁重工作；之后所有的访问都直接取回已有的结果。

现在假设我还想在这个数据库系统中加入事务（transaction）。下次用户访问某个动态属性时，我想知道数据库中对应的记录是否仍然有效，以及事务是否仍然处于打开状态。\texttt{\_\_getattr\_\_} 钩子并不会在每次访问属性时都被调用，因为它会把对象的实例字典作为已有属性的快速通道（fast path）。

为了实现这种更高级的用例，Python 还提供了另一个 \texttt{object} 钩子，名为 \texttt{\_\_getattribute\_\_}。这个特殊方法在每次访问对象上的属性时都会被调用，即便该属性确实存在于属性字典中也不例外。这让我可以在每次属性访问时做一些事情，比如检查全局事务状态。需要注意的是，这样的操作可能会带来显著的开销（overhead）并对性能产生负面影响，但有时这么做是值得的。这里我定义 \texttt{ValidatingRecord}，以便在每次调用 \texttt{\_\_getattribute\_\_} 时记录日志：

\begin{codebox}[python]
class ValidatingRecord:
    def __init__(self):
        self.exists = 5
    def __getattribute__(self, name):
        print(f"* Called __getattribute__({name!r})")
        try:
            value = super().__getattribute__(name)
            print(f"* Found {name!r}, returning {value!r}")
            return value
        except AttributeError:
            value = f"Value for {name}"
            print(f"* Setting {name!r} to {value!r}")
            setattr(self, name, value)
            return value
data = ValidatingRecord()
print("exists:     ", data.exists)
print("First foo:  ", data.foo)
print("Second foo: ", data.foo)
>>>
* Called __getattribute__('exists')
* Found 'exists', returning 5
exists:      5
* Called __getattribute__('foo')
* Setting 'foo' to 'Value for foo'
First foo:   Value for foo
* Called __getattribute__('foo')
* Found 'foo', returning 'Value for foo'
Second foo:  Value for foo
\end{codebox}

如果某个被动态访问的属性本不应该存在，你可以抛出 \texttt{AttributeError}，从而让 \texttt{\_\_getattr\_\_} 和 \texttt{\_\_getattribute\_\_} 都触发 Python 标准的"属性缺失"行为：

\begin{codebox}[python]
class MissingPropertyRecord:
    def __getattr__(self, name):
        if name == "bad_name":
            raise AttributeError(f"{name} is missing")
        ...
data = MissingPropertyRecord()
data.bad_name
>>>
Traceback ...
AttributeError: bad_name is missing
\end{codebox}

实现通用功能的 Python 代码常常依赖 \texttt{hasattr} 这个内置（built-in）函数来判断属性是否存在，并使用 \texttt{getattr} 这个内置函数来取回属性值。这两个函数在调用 \texttt{\_\_getattr\_\_} 之前，也会先在实例字典里查找属性名：

\begin{codebox}[python]
data = LoggingLazyRecord()  # Implements __getattr__
print("Before:         ", data.__dict__)
print("Has first foo:  ", hasattr(data, "foo"))
print("After:          ", data.__dict__)
print("Has second foo: ", hasattr(data, "foo"))
>>>
Before:          {'exists': 5}
* Called __getattr__('foo'), populating instance dictionary
* Returning 'Value for foo'
Has first foo:   True
After:           {'exists': 5, 'foo': 'Value for foo'}
Has second foo:  True
\end{codebox}

在上例中，\texttt{\_\_getattr\_\_} 只被调用了一次（对应第一次 \texttt{hasattr} 调用）。相比之下，实现了 \texttt{\_\_getattribute\_\_} 的类，每当对实例使用 \texttt{hasattr} 或 \texttt{getattr} 时，都会调用该方法：

\begin{codebox}[python]
data = ValidatingRecord()  # Implements __getattribute__
print("Has first foo:  ", hasattr(data, "foo"))
print("Has second foo: ", hasattr(data, "foo"))
>>>
* Called __getattribute__('foo')
* Setting 'foo' to 'Value for foo'
Has first foo:   True
* Called __getattribute__('foo')
* Found 'foo', returning 'Value for foo'
Has second foo:  True
\end{codebox}

现在，假设我想在把值赋给我的 Python 对象时，惰性地把数据推送回数据库。我可以用 \texttt{\_\_setattr\_\_} 来实现，它是 \texttt{object} 的一个类似钩子，让你可以拦截任意的属性赋值。与使用 \texttt{\_\_getattr\_\_} 和 \texttt{\_\_getattribute\_\_} 取回属性不同，这里不需要两个独立的方法。\texttt{\_\_setattr\_\_} 方法在每次给实例上的属性赋值时（无论是直接赋值，还是通过 \texttt{setattr} 内置函数）都会被调用：

\begin{codebox}[python]
class SavingRecord:
    def __setattr__(self, name, value):
        # 为记录保存一些数据
        ...
        super().__setattr__(name, value)
\end{codebox}

这里我定义一个 \texttt{SavingRecord} 的日志子类。它的 \texttt{\_\_setattr\_\_} 方法在每次属性赋值时都会被调用：

\begin{codebox}[python]
class LoggingSavingRecord(SavingRecord):
    def __setattr__(self, name, value):
        print(f"* Called __setattr__({name!r}, {value!r})")
        super().__setattr__(name, value)
data = LoggingSavingRecord()
print("Before: ", data.__dict__)
data.foo = 5
print("After:  ", data.__dict__)
data.foo = 7
print("Finally:", data.__dict__)
>>>
Before:  {}
* Called __setattr__('foo', 5)
After:   {'foo': 5}
* Called __setattr__('foo', 7)
Finally: {'foo': 7}
\end{codebox}

\texttt{\_\_getattribute\_\_} 和 \texttt{\_\_setattr\_\_} 的问题在于，它们会在对象的每次属性访问时都被调用——即便你并不想这样。例如，假设访问我的对象上的属性，实际上应该去查找某个关联字典里的键：

\begin{codebox}[python]
class BrokenDictionaryRecord:
    def __init__(self, data):
        self._data = data
    def __getattribute__(self, name):
        print(f"* Called __getattribute__({name!r})")
        return self._data[name]
\end{codebox}

而这需要从 \texttt{\_\_getattribute\_\_} 方法里访问 \texttt{self.\_data}。然而，如果我真的这样做，Python 会一直递归下去，直到达到栈（stack）上限，随后程序就会崩溃：

\begin{codebox}[python]
data = BrokenDictionaryRecord({"foo": 3})
data.foo
>>>
* Called __getattribute__('foo')
* Called __getattribute__('_data')
* Called __getattribute__('_data')
* Called __getattribute__('_data')
...
Traceback ...
RecursionError: maximum recursion depth exceeded while calling
a Python object
\end{codebox}

问题在于，\texttt{\_\_getattribute\_\_} 会访问 \texttt{self.\_data}，而这又会导致 \texttt{\_\_getattribute\_\_} 再次运行，进而再次访问 \texttt{self.\_data}，如此循环往复。解决办法是使用 \texttt{super().\_\_getattribute\_\_} 方法从实例属性字典中取值；这样可以避免意外的递归：

\begin{codebox}[python]
class DictionaryRecord:
    def __init__(self, data):
        self._data = data
    def __getattribute__(self, name):
        print(f"* Called __getattribute__({name!r})")
        data_dict = super().__getattribute__("_data")
        return data_dict[name]
data = DictionaryRecord({"foo": 3})
print("foo: ", data.foo)
>>>
* Called __getattribute__('foo')
foo:  3
\end{codebox}

与之类似，那些会修改对象上属性的 \texttt{\_\_setattr\_\_} 方法，也需要使用 \texttt{super().\_\_setattr\_\_}。

\subsection{需要记住的要点}

\begin{itemize}
\item 使用 \texttt{\_\_getattr\_\_} 和 \texttt{\_\_setattr\_\_} 来惰性加载并保存对象的属性。
\item 要理解，\texttt{\_\_getattr\_\_} 只在访问缺失的属性时才会被调用，而 \texttt{\_\_getattribute\_\_} 在每次访问任何属性时都会被调用。
\item 在 \texttt{\_\_getattribute\_\_} 和 \texttt{\_\_setattr\_\_} 方法的实现中，通过调用 \texttt{super().\_\_getattribute\_\_} 和 \texttt{super().\_\_getattr\_\_} 来访问对象属性，从而避免无限递归。
\end{itemize}

\section{Item 62: 用 \texttt{\_\_init\_subclass\_\_} 校验子类}

元类（metaclass）最简单的应用之一，就是校验一个类是否被正确定义。当你在构建复杂的类层次结构（class hierarchy）时，可能想强制某种编码风格、要求重写某些方法，或者在类的属性之间建立严格的关系。元类通过提供一种可靠的方式，让验证代码在每次定义新子类时都能运行，从而满足了这些需求。

类的验证代码通常运行在 \texttt{\_\_init\_\_} 方法中，在运行时构造该类型的对象时进行（参见 Item 58（"用普通属性代替 setter 和 getter 方法"）的例子）。使用元类来做验证，可以让错误更早地暴露出来，例如在程序启动时、包含该类的模块首次被导入时就会报错。

在介绍如何定义用于校验子类的元类之前，先理解元类对标准对象做了什么很重要。元类通过继承 \texttt{type} 来定义。一个类通过在继承参数列表中传入 \texttt{metaclass} 关键字参数来指定其元类。通常，当任何相关的 \texttt{class} 语句出现时，元类的 \texttt{\_\_new\_\_} 方法会被调用，并接收到该语句的内容。下面我用一个基础元类，在实际构造类型之前检查类的信息：

\begin{codebox}[python]
class Meta(type):
    def __new__(meta, name, bases, class_dict):
        print(f"* Running {meta}.__new__ for {name}")
        print("Bases:", bases)
        print(class_dict)
        return type.__new__(meta, name, bases, class_dict)
class MyClass(metaclass=Meta):
    stuff = 123
    def foo(self):
        pass
class MySubclass(MyClass):
    other = 567
    def bar(self):
        pass
>>>
* Running <class '__main__.Meta'>.__new__ for MyClass
Bases: ()
{'__module__': '__main__',
 '__qualname__': 'MyClass',
 'stuff': 123,
 'foo': <function MyClass.foo at 0x104a63a60>}
* Running <class '__main__.Meta'>.__new__ for MySubclass
Bases: (<class '__main__.MyClass'>,)
{'__module__': '__main__',
 '__qualname__': 'MySubclass',
 'other': 567,
 'bar': <function MySubclass.bar at 0x104a63b00>}
\end{codebox}

元类可以访问类的名称、它所继承的父类（\texttt{bases}），以及在该类体中定义的所有类属性。所有类都继承自 \texttt{object}，所以它不会显式地出现在基类的元组中。

我可以给 \texttt{Meta.\_\_new\_\_} 方法添加功能，以便在相关子类被定义之前校验其所有参数。例如，假设我想表示任意类型的多边形（polygon）。我可以定义一个特殊的校验元类，并在多边形类层次的基类中加以使用。注意，重要的是不要把同样的校验应用到基类本身：

\begin{codebox}[python]
class ValidatePolygon(type):
    def __new__(meta, name, bases, class_dict):
        # 只校验 Polygon 类的子类
        if bases:
            if class_dict["sides"] < 3:
                raise ValueError("Polygons need 3+ sides")
        return type.__new__(meta, name, bases, class_dict)
class Polygon(metaclass=ValidatePolygon):
    sides = None  # 必须由子类指定
    @classmethod
    def interior_angles(cls):
        return (cls.sides - 2) * 180
class Triangle(Polygon):
    sides = 3
class Rectangle(Polygon):
    sides = 4
class Nonagon(Polygon):
    sides = 9
assert Triangle.interior_angles() == 180
assert Rectangle.interior_angles() == 360
assert Nonagon.interior_angles() == 1260
\end{codebox}

如果我试图定义一个边数少于 3 的多边形，校验逻辑会让 \texttt{class} 语句在类体执行完毕后立即失败。这意味着，当我定义这样一个类时，程序甚至根本无法开始运行（除非它被定义在动态导入的模块中；参见 Item 98（"用动态导入惰性加载模块以减少启动时间"）了解这是如何发生的）：

\begin{codebox}[python]
print("Before class")
class Line(Polygon):
    print("Before sides")
    sides = 2
    print("After sides")
print("After class")
>>>
Before class
Before sides
After sides
Traceback ...
ValueError: Polygons need 3+ sides
\end{codebox}

为了让 Python 完成这样一个基础的任务，这看起来需要相当多的机制。幸运的是，Python 3.6 引入了简化语法——\texttt{\_\_init\_subclass\_\_} 这个特殊的类方法——来实现同样的行为，并且完全避免了元类。下面我用这一机制提供与之前同等级的校验：

\begin{codebox}[python]
class BetterPolygon:
    sides = None  # 必须由子类指定
    def __init_subclass__(cls):
        super().__init_subclass__()
        if cls.sides < 3:
            raise ValueError("Polygons need 3+ sides")
    @classmethod
    def interior_angles(cls):
        return (cls.sides - 2) * 180
class Hexagon(BetterPolygon):
    sides = 6
assert Hexagon.interior_angles() == 720
\end{codebox}

现在代码短了许多，而且 \texttt{ValidatePolygon} 元类彻底消失了。可读性也更好了，因为我可以直接在 \texttt{cls} 实例上访问 \texttt{sides} 属性，而不必通过 \texttt{class\_dict["sides"]} 深入到类的字典中。

如果我定义一个 \texttt{BetterPolygon} 的无效子类，就会抛出与之前相同的异常：

\begin{codebox}[python]
print("Before class")
class Point(BetterPolygon):
    sides = 1
print("After class")
>>>
Before class
Traceback ...
ValueError: Polygons need 3+ sides
\end{codebox}

标准 Python 元类机制的另一个问题是，每个类定义只能指定一个元类。下面我定义第二个元类，想用它来校验某个区域所用填充颜色（不一定只是多边形）：

\begin{codebox}[python]
class ValidateFilled(type):
    def __new__(meta, name, bases, class_dict):
        # 只校验 Filled 类的子类
        if bases:
            if class_dict["color"] not in ("red", "green"):
                raise ValueError("Fill color must be supported")
        return type.__new__(meta, name, bases, class_dict)
class Filled(metaclass=ValidateFilled):
    color = None  # 必须由子类指定
\end{codebox}

当我试图同时使用 \texttt{Polygon} 元类和 \texttt{Filled} 元类时，会得到一条含糊的错误信息：

\begin{codebox}[python]
class RedPentagon(Filled, Polygon):
    color = "blue"
    sides = 5
>>>
Traceback ...
TypeError: metaclass conflict: the metaclass of a derived class
'must be a (non-strict) subclass of the metaclasses of all its
'bases
\end{codebox}

可以通过创建复杂的元类 \texttt{type} 定义层次来叠加校验，从而修复这个问题：

\begin{codebox}[python]
class ValidatePolygon(type):
    def __new__(meta, name, bases, class_dict):
        # 只校验非根类
        if not class_dict.get("is_root"):
            if class_dict["sides"] < 3:
                raise ValueError("Polygons need 3+ sides")
        return type.__new__(meta, name, bases, class_dict)
class Polygon(metaclass=ValidatePolygon):
    is_root = True
    sides = None  # 必须由子类指定
class ValidateFilledPolygon(ValidatePolygon):
    def __new__(meta, name, bases, class_dict):
        # 只校验非根类
        if not class_dict.get("is_root"):
            if class_dict["color"] not in ("red", "green"):
                raise ValueError("Fill color must be supported")
        return super().__new__(meta, name, bases, class_dict)
class FilledPolygon(Polygon, metaclass=ValidateFilledPolygon):
    is_root = True
    color = None  # 必须由子类指定
\end{codebox}

这要求每一个 \texttt{FilledPolygon} 实例都必须是 \texttt{Polygon} 的实例：

\begin{codebox}[python]
class GreenPentagon(FilledPolygon):
    color = "green"
    sides = 5
greenie = GreenPentagon()
assert isinstance(greenie, Polygon)
\end{codebox}

对颜色的校验同样有效：

\begin{codebox}[python]
class OrangePentagon(FilledPolygon):
    color = "orange"
    sides = 5
>>>
Traceback ...
ValueError: Fill color must be supported
\end{codebox}

对边的数量校验同样有效：

\begin{codebox}[python]
class RedLine(FilledPolygon):
    color = "red"
    sides = 2
>>>
Traceback ...
ValueError: Polygons need 3+ sides
\end{codebox}

然而，这种做法破坏了可组合性（composability），而这通常正是此类类校验的目的（类似于 mix-in；参见 Item 54（"用 mix-in 类组合功能"））。如果我想把 \texttt{ValidateFilledPolygon} 中的颜色校验逻辑应用到另一组类层次上，就不得不再次复制全部逻辑，这降低了代码复用并增加了样板代码（boilerplate）。

\texttt{\_\_init\_subclass\_\_} 这个特殊的类方法也可以用来解决这个问题。只要使用 \texttt{super} 内置函数来调用任何父级或兄弟级的 \texttt{\_\_init\_subclass\_\_} 定义，它就可以在类层次的多个层级上被定义。下面我定义一个表示区域填充颜色的类，可以跟前面的 \texttt{BetterPolygon} 类组合：

\begin{codebox}[python]
class Filled:
    color = None  # 必须由子类指定
    def __init_subclass__(cls):
        super().__init_subclass__()
        if cls.color not in ("red", "green", "blue"):
            raise ValueError("Fills need a valid color")
\end{codebox}

我可以从这两个类继承来定义新类。两个类都会调用 \texttt{super().\_\_init\_subclass\_\_()}，使得它们各自的校验逻辑在子类被创建时运行：

\begin{codebox}[python]
class RedTriangle(Filled, BetterPolygon):
    color = "red"
    sides = 3
ruddy = RedTriangle()
assert isinstance(ruddy, Filled)
assert isinstance(ruddy, BetterPolygon)
\end{codebox}

如果我边数指定错误，就会得到一个校验错误：

\begin{codebox}[python]
print("Before class")
class BlueLine(Filled, BetterPolygon):
    color = "blue"
    sides = 2
print("After class")
>>>
Before class
Traceback ...
ValueError: Polygons need 3+ sides
\end{codebox}

如果我颜色指定错误，同样会得到一个校验错误：

\begin{codebox}[python]
print("Before class")
class BeigeSquare(Filled, BetterPolygon):
    color = "beige"
    sides = 4
print("After class")
>>>
Before class
Traceback ...
ValueError: Fills need a valid color
\end{codebox}

你甚至可以在多重继承（multiple inheritance）、菱形继承（diamond inheritance）等复杂情况下使用 \texttt{\_\_init\_subclass\_\_}（背景知识参见 Item 53（"用 super 初始化父类"））。下面我定义一个基础的菱形层次来演示它的实际运作：

\begin{codebox}[python]
class Top:
    def __init_subclass__(cls):
        super().__init_subclass__()
        print(f"Top for {cls}")
class Left(Top):
    def __init_subclass__(cls):
        super().__init_subclass__()
        print(f"Left for {cls}")
class Right(Top):
    def __init_subclass__(cls):
        super().__init_subclass__()
        print(f"Right for {cls}")
class Bottom(Left, Right):
    def __init_subclass__(cls):
        super().__init_subclass__()
        print(f"Bottom for {cls}")
>>>
Top for <class '__main__.Left'>
Top for <class '__main__.Right'>
Top for <class '__main__.Bottom'>
Right for <class '__main__.Bottom'>
Left for <class '__main__.Bottom'>
\end{codebox}

正如预期，\texttt{Top.\_\_init\_subclass\_\_} 对每个类只被调用一次，即使对于 \texttt{Bottom} 类而言，通过它的 \texttt{Left} 和 \texttt{Right} 父类有两条路径可以到达它。

\subsection{需要记住的要点}

\begin{itemize}
\item 元类的 \texttt{\_\_new\_\_} 方法会在 \texttt{class} 语句的整个类体处理完毕后才运行。
\item 元类可以用来在类被定义之后、创建之前检查或修改它，但它们往往比你所需要的更"重"（heavyweight）。
\item 使用 \texttt{\_\_init\_subclass\_\_} 来确保在子类被定义的时刻（即该类型的对象被构造之前）其形式是良好的。
\item 务必从你的类的 \texttt{\_\_init\_subclass\_\_} 定义中调用 \texttt{super().\_\_init\_subclass\_\_()}，以便在类的多个层次和多重继承中实现可组合的校验。
\end{itemize}

\section{Item 63: 用 \texttt{\_\_init\_subclass\_\_} 注册类的存在}

元类（metaclass）的另一个常见用途（背景参见 Item 62：用 \texttt{\_\_init\_subclass\_\_} 校验子类）是在程序中自动注册类型。注册对于执行反向查找（reverse lookup）很有用，在这种情况下，你需要把一个标识符（identifier）映射回相应的类（class）。

例如，假设我想用 JSON 实现自己的 Python 对象序列化表示。我需要一种把对象转换成 JSON 字符串的方法。在这里，我是通过定义一个基类（base class）来通用地做到这一点的：该基类记录构造函数（constructor）的参数，并把它们变成一个 JSON 字典（dictionary）（另一种做法参见 Item 54：用 mix-in（混入类）组合功能）：

\begin{codebox}[python]
import json
class Serializable:
    def __init__(self, *args):
        self.args = args
    def serialize(self):
        return json.dumps({"args": self.args})
\end{codebox}

这个类让把简单数据结构序列化成字符串变得很容易，像下面这样：

\begin{codebox}[python]
class Point2D(Serializable):
    def __init__(self, x, y):
        super().__init__(x, y)
        self.x = x
        self.y = y
    def __repr__(self):
        return f"Point2D({self.x}, {self.y})"
point = Point2D(5, 3)
print("Object:    ", point)
print("Serialized:", point.serialize())
>>>
Object:     Point2D(5, 3)
Serialized: {"args": [5, 3]}
\end{codebox}

现在我需要反序列化（deserialize）这个 JSON 字符串，并构造它所对应的 \texttt{Point2D} 对象（object）。在这里，我定义另一个类，它可以从其 \texttt{Serializable} 父类（parent class）反序列化数据（背景参见 Item 52：用 \texttt{@classmethod} 多态（polymorphism）通用地构造对象）：

\begin{codebox}[python]
class Deserializable(Serializable):
    @classmethod
    def deserialize(cls, json_data):
        params = json.loads(json_data)
        return cls(*params["args"])
\end{codebox}

把 \texttt{Deserializable} 作为父类，可以很方便地以通用方式对简单对象进行序列化和反序列化：

\begin{codebox}[python]
class BetterPoint2D(Deserializable):
    def __init__(self, x, y):
        super().__init__(x, y)
        self.x = x
        self.y = y
    def __repr__(self):
        return f"Point2D({self.x}, {self.y})"
before = BetterPoint2D(5, 3)
print("Before:    ", before)
data = before.serialize()
print("Serialized:", data)
after = BetterPoint2D.deserialize(data)
print("After:     ", after)
>>>
Before:     Point2D(5, 3)
Serialized: {"args": [5, 3]}
After:      Point2D(5, 3)
\end{codebox}

这种做法的问题在于，只有当你在事先就知道序列化数据的预期类型时，它才有效（例如 \texttt{Point2D}、\texttt{BetterPoint2D}）。理想情况下，你会有大量类序列化到 JSON，并且有一个通用函数可以把它们中的任意一个反序列化回相应的 Python 对象（object）（类似的例子见 Item 50：用 \texttt{functools.singledispatch} 做函数式编程，而非面向对象的多态）。

要做到这一点，我可以把序列化对象的类名放进 JSON 数据里：

\begin{codebox}[python]
class BetterSerializable:
    def __init__(self, *args):
        self.args = args
    def serialize(self):
        return json.dumps(
            {
                "class": self.__class__.__name__,
                "args": self.args,
            }
        )
    def __repr__(self):
        name = self.__class__.__name__
        args_str = ", ".join(str(x) for x in self.args)
        return f"{name}({args_str})"
\end{codebox}

然后我可以维护一张从类名映射回这些对象的构造函数（constructor）的映射表（mapping）。这个通用的 \texttt{deserialize} 函数适用于传给 \texttt{register\_class} 的任何类：

\begin{codebox}[python]
REGISTRY = {}
def register_class(target_class):
    REGISTRY[target_class.__name__] = target_class
def deserialize(data):
    params = json.loads(data)
    name = params["class"]
    target_class = REGISTRY[name]
    return target_class(*params["args"])
\end{codebox}

为了确保 \texttt{deserialize} 始终正常工作，我必须为每一个将来可能要反序列化的类调用 \texttt{register\_class}：

\begin{codebox}[python]
class EvenBetterPoint2D(BetterSerializable):
    def __init__(self, x, y):
        super().__init__(x, y)
        self.x = x
        self.y = y
register_class(EvenBetterPoint2D)
\end{codebox}

现在我可以对任意 JSON 字符串进行反序列化，而不必先知道它包含哪个类：

\begin{codebox}[python]
before = EvenBetterPoint2D(5, 3)
print("Before:    ", before)
data = before.serialize()
print("Serialized:", data)
after = deserialize(data)
print("After:     ", after)
>>>
Before:     EvenBetterPoint2D(5, 3)
Serialized: {"class": "EvenBetterPoint2D", "args": [5, 3]}
After:      EvenBetterPoint2D(5, 3)
\end{codebox}

这种做法的问题在于，有可能忘记调用 \texttt{register\_class}：

\begin{codebox}[python]
class Point3D(BetterSerializable):
    def __init__(self, x, y, z):
        super().__init__(x, y, z)
        self.x = x
        self.y = y
        self.z = z
# 忘记调用 register_class 了！哎呀！
\end{codebox}

当我试图反序列化一个我忘记注册的类的对象（object）时，代码会在运行时崩溃：

\begin{codebox}[python]
point = Point3D(5, 9, -4)
data = point.serialize()
deserialize(data)
>>>
Traceback ...
KeyError: 'Point3D'
\end{codebox}

即使我选择了继承 \texttt{BetterSerializable}，如果我在 \texttt{class} 语句（statement）体之后忘记调用 \texttt{register\_class}，实际上也无法获得它的全部特性。这种做法容易出错，而且对新手来说尤其难以调试。类装饰器（class decorator）也会出现同样的疏漏（参见 Item 66：在需要可组合类扩展时，优先使用类装饰器而非元类）。

如果我能对程序员使用 \texttt{BetterSerializable} 的意图做出响应，并确保在所有情况下都调用 \texttt{register\_class}，那会怎样？元类（metaclass）可以通过在定义子类（subclass）时拦截 \texttt{class} 语句来实现这一点。在这里，我用一个元类（metaclass）和相应的父类（superclass），在子类的 \texttt{class} 语句结束后立即注册它们：

\begin{codebox}[python]
class Meta(type):
    def __new__(meta, name, bases, class_dict):
        cls = type.__new__(meta, name, bases, class_dict)
        register_class(cls)
        return cls
class RegisteredSerializable(BetterSerializable, metaclass=Meta):
    pass
\end{codebox}

当我定义 \texttt{RegisteredSerializable} 的子类（subclass）时，我可以确信对 \texttt{register\_class} 的调用会发生，并且 \texttt{deserialize} 将始终按预期工作：

\begin{codebox}[python]
class Vector3D(RegisteredSerializable):
    def __init__(self, x, y, z):
        super().__init__(x, y, z)
        self.x, self.y, self.z = x, y, z
before = Vector3D(10, -7, 3)
print("Before:    ", before)
data = before.serialize()
print("Serialized:", data)
print("After:     ", deserialize(data))
>>>
Before:     Vector3D(10, -7, 3)
Serialized: {"class": "Vector3D", "args": [10, -7, 3]}
After:      Vector3D(10, -7, 3)
\end{codebox}

更好的做法是使用 \texttt{\_\_init\_subclass\_\_} 这个特殊的类方法（class method）。这种简化的语法（syntax）自 Python 3.6 引入，减少了在定义类时应用自定义逻辑（logic）所产生的视觉干扰。它也让那些可能被元类（metaclass）语法的复杂性（complexity）搞糊涂的新手更容易上手。在这里，我实现一个新的父类（superclass），让它自动调用 \texttt{register\_class}，以及一个使用它的子类（subclass）：

\begin{codebox}[python]
class BetterRegisteredSerializable(BetterSerializable):
    def __init_subclass__(cls):
        super().__init_subclass__()
        register_class(cls)
class Vector1D(BetterRegisteredSerializable):
    def __init__(self, magnitude):
        super().__init__(magnitude)
        self.magnitude = magnitude
\end{codebox}

序列化与反序列化对这个新类都如期工作：

\begin{codebox}[python]
before = Vector1D(6)
print("Before:    ", before)
data = before.serialize()
print("Serialized:", data)
print("After:     ", deserialize(data))
>>>
Before:     Vector1D(6)
Serialized: {"class": "Vector1D", "args": [6]}
After:      Vector1D(6)
\end{codebox}

通过对类注册使用 \texttt{\_\_init\_subclass\_\_}（或元类（metaclass）），只要继承树（inheritance tree）是正确的，你就能确保永远不会漏掉注册某个类。正如我所展示的，这很适用于序列化（serialization），也适用于数据库的对象关系映射（ORM）、可扩展的插件（plug-in）系统，以及回调钩子（callback hook）。

\subsection{需要记住的要点}

\begin{itemize}
\item 类注册是构建模块化（modular）Python 程序的一种有用模式（pattern）。
\item 元类让你在程序中每次有基类（base class）被子类化（subclass）时，自动运行注册代码。
\item 通过确保你永远不会漏掉一次注册调用，用元类做类注册可以帮你避免错误（error）。
\item 优先使用 \texttt{\_\_init\_subclass\_\_} 而非标准的元类机制（metaclass machinery），因为它更清晰，也更容易被新手理解。
\end{itemize}

\section{Item 64: 用 \texttt{\_\_set\_name\_\_} 标注类属性}

元类（metaclass）还带来了一个实用特性（背景参见 Item 62（"用 \texttt{\_\_init\_subclass\_\_} 校验子类"））：在类定义完成之后、实际使用之前，具备修改或标注其属性的能力。这种做法通常与描述符（descriptor，细节参见 Item 60（"为可复用的 \texttt{@property} 方法使用描述符"））搭配使用，让这些属性能够更深入地自省（introspection）它们在所属类中的使用方式。

举例来说，假设我想定义一个新类来表示客户数据库中的一行。我希望为数据库表的每一列，在类上都有一个对应的属性。下面我定义一个描述符类，把属性与列名关联起来：

\begin{codebox}[python]
class Field:
    def __init__(self, column_name):
        self.column_name = column_name
        self.internal_name = "_" + self.column_name
\end{codebox}

借助内置（built-in）函数 \texttt{setattr}，我可以用列名把所有逐实例的状态以受保护字段的形式直接保存到实例字典（instance dictionary）里；之后，可以用 \texttt{getattr} 来加载状态（背景参见 Item 61（"用 \texttt{\_\_getattr\_\_}、\texttt{\_\_getattribute\_\_} 和 \texttt{\_\_setattr\_\_} 实现惰性属性"））：

\begin{codebox}[python]
    def __get__(self, instance, instance_type):
        if instance is None:
            return self
        return getattr(instance, self.internal_name, "")
    def __set__(self, instance, value):
        setattr(instance, self.internal_name, value)
\end{codebox}

定义表示一行的类时，需要为每个描述符属性提供数据库表的列名：

\begin{codebox}[python]
class Customer:
    # Class attributes
    first_name = Field("first_name")
    last_name = Field("last_name")
    prefix = Field("prefix")
    suffix = Field("suffix")
\end{codebox}

使用这个行类很简单。这里，\texttt{Field} 描述符会按预期修改实例字典 \texttt{\_\_dict\_\_}：

\begin{codebox}[python]
cust = Customer()
print(f"Before: {cust.first_name!r} {cust.__dict__}")
cust.first_name = "Euclid"
print(f"After:  {cust.first_name!r} {cust.__dict__}")
>>>
Before: '' {}
After:  'Euclid' {'_first_name': 'Euclid'}
\end{codebox}

但这个类定义的代码似乎存在冗余。我已经在左侧为类声明了字段名（\texttt{field\_name =}）。为什么我还必须把包含相同信息的字符串传给右侧的 \texttt{Field} 构造函数（\texttt{Field("first\_name")}）？

\begin{codebox}[python]
class Customer:
    # 左侧与右侧重复
    first_name = Field("first_name")
    ...
\end{codebox}

问题在于，\texttt{Customer} 类定义的计算顺序，与从左往右阅读的顺序正好相反。首先，\texttt{Field} 构造函数以 \texttt{Field("first\_name")} 的形式被调用；然后，其返回值被赋给 \texttt{Customer.first\_name} 类属性（class attribute）。\texttt{Field} 实例无法提前知道它会被赋给哪个类属性。

为了消除这种冗余，我可以使用元类。元类让你能够直接钩入（hook）\texttt{class} 语句，并在类体（class body）一结束就采取行动。在本例中，我可以用元类自动为描述符赋值 \texttt{Field.column\_name} 和 \texttt{Field.internal\_name}，而无需多次手动指定字段名：

\begin{codebox}[python]
class Meta(type):
    def __new__(meta, name, bases, class_dict):
        for key, value in class_dict.items():
            if isinstance(value, Field):
                value.column_name = key
                value.internal_name = "_" + key
        cls = type.__new__(meta, name, bases, class_dict)
        return cls
\end{codebox}

这里我定义一个使用元类的基类（base class）。所有表示数据库行的类都应继承自这个类，以确保它们使用元类：

\begin{codebox}[python]
class DatabaseRow(metaclass=Meta):
    pass
\end{codebox}

为了配合元类，字段描述符基本不变。唯一的不同是它不再需要向其构造函数传入任何参数；相反，它的属性由上面的 \texttt{Meta.\_\_new\_\_} 方法设置：

\begin{codebox}[python]
class Field:
    def __init__(self):
        # 这些将由元类赋值。
        self.column_name = None
        self.internal_name = None
    def __get__(self, instance, instance_type):
        if instance is None:
            return self
        return getattr(instance, self.internal_name, "")
    def __set__(self, instance, value):
        setattr(instance, self.internal_name, value)
\end{codebox}

当我使用元类、新的 \texttt{DatabaseRow} 基类和新的 \texttt{Field} 描述符时，数据库行的类定义不再有之前的冗余：

\begin{codebox}[python]
class BetterCustomer(DatabaseRow):
    first_name = Field()
    last_name = Field()
    prefix = Field()
    suffix = Field()
\end{codebox}

新类的行为与旧类完全相同：

\begin{codebox}[python]
cust = BetterCustomer()
print(f"Before: {cust.first_name!r} {cust.__dict__}")
cust.first_name = "Euler"
print(f"After:  {cust.first_name!r} {cust.__dict__}")
>>>
Before: '' {}
After:  'Euler' {'_first_name': 'Euler'}
\end{codebox}

这种做法的麻烦在于，除非同时继承自 \texttt{DatabaseRow}，否则无法把 \texttt{Field} 类用于属性。如果你不小心忘了继承 \texttt{DatabaseRow}，或者因为类层次结构的其他结构性要求而不想继承，代码就会出错：

\begin{codebox}[python]
class BrokenCustomer:  # 缺少继承
    first_name = Field()
    last_name = Field()
    prefix = Field()
    suffix = Field()
cust = BrokenCustomer()
cust.first_name = "Mersenne"
>>>
Traceback ...
TypeError: attribute name must be string, not 'NoneType'
\end{codebox}

解决这个问题的办法，是为描述符使用 \texttt{\_\_set\_name\_\_} 特殊方法。这个方法在 Python 3.6 中引入，会在包含它的类被定义时，对每一个描述符实例调用。它接收两个参数：包含该描述符实例的所属类（owning class），以及该描述符实例被赋给的属性名。这里我完全不定义元类，而是把上面 \texttt{Meta.\_\_new\_\_} 方法所做的事情移到了 \texttt{\_\_set\_name\_\_} 方法里：

\begin{codebox}[python]
class Field:
    def __init__(self):
        self.column_name = None
        self.internal_name = None
    def __set_name__(self, owner, column_name):
        # 类创建时为每个描述符调用
        self.column_name = column_name
        self.internal_name = "_" + column_name
    def __get__(self, instance, instance_type):
        if instance is None:
            return self
        return getattr(instance, self.internal_name, "")
    def __set__(self, instance, value):
        setattr(instance, self.internal_name, value)
\end{codebox}

现在，我既能享受 \texttt{Field} 描述符的好处，又不必继承自某个特定的父类，也不必使用元类：

\begin{codebox}[python]
class FixedCustomer:  # 没有父类
    first_name = Field()
    last_name = Field()
    prefix = Field()
    suffix = Field()
cust = FixedCustomer()
print(f"Before: {cust.first_name!r} {cust.__dict__}")
cust.first_name = "Mersenne"
print(f"After:  {cust.first_name!r} {cust.__dict__}")
>>>
Before: '' {}
After:  'Mersenne' {'_first_name': 'Mersenne'}
\end{codebox}

\subsection{需要记住的要点}

\begin{itemize}
\item 元类让你能够在类完全定义之前修改它的属性。
\item 描述符与元类组合在一起，能为实现声明式行为和运行时自省（runtime introspection）提供强大的能力。
\item 在你的描述符类上定义 \texttt{\_\_set\_name\_\_}，让它们能够考虑其所在的类以及它的属性名。
\end{itemize}

\section{Item 65: 考虑类体定义顺序以确立属性间的关系}

Python 程序中定义的许多类的用途，是表示在别处创建和维护的外部数据。例如，假设我有一个 CSV（逗号分隔值，comma-separated values）文件，其中包含一份货运清单，每一行都包括目的地城市、运输方式和货物重量。这里我使用 \texttt{csv} 内置（built-in）模块读入这些数据：

\begin{codebox}[python]
import csv
with open("packages.csv") as f:
    for row in csv.reader(f):
        print(row)
>>>
['Sydney', 'truck', '25']
['Melbourne', 'boat', '6']
['Brisbane', 'plane', '12']
['Perth', 'road train', '90']
['Adelaide', 'truck', '17']
...
\end{codebox}

我可以定义一个新类来存储这些数据，并定义一个辅助函数，给定一个 CSV 行就能创建出对象（背景知识见 Item 52（「使用 \texttt{@classmethod} 多态以通用方式构造对象」））：

\begin{codebox}[python]
class Delivery:
    def __init__(self, destination, method, weight):
        self.destination = destination
        self.method = method
        self.weight = weight
    @classmethod
    def from_row(cls, row):
        return cls(row[0], row[1], row[2])
\end{codebox}

当提供一列值（每列一个）时，它的运行符合预期：

\begin{codebox}[python]
row1 = ["Sydney", "truck", "25"]
obj1 = Delivery.from_row(row1)
print(obj1.__dict__)
>>>
{'destination': 'Sydney', 'method': 'truck', 'weight': '25'}
\end{codebox}

如果 CSV 文件增加了更多列，或者列的顺序被调整，我只需稍作努力，就能对 \texttt{\_\_init\_\_} 和 \texttt{from\_row} 方法做相应调整，以维持与文件格式的兼容。现在设想我想处理很多种 CSV 文件，每种都有不同数量的列和不同类型的单元格取值。如果我能以更少的样板代码，为每种 CSV 文件高效地定义一个新类，那就更好了。

这里我尝试通过一个基类来实现这一点：它使用 \texttt{fields} 类属性（class attribute）把 CSV 列（按它们在文件中出现的顺序）映射到对象属性名（另一种做法见 Item 64（「用 \texttt{\_\_set\_name\_\_} 标注类属性」））：

\begin{codebox}[python]
class RowMapper:
    fields = ()  # 必须按 CSV 列顺序排列
    def __init__(self, **kwargs):
        for key, value in kwargs.items():
            if key not in type(self).fields:
                raise TypeError(f"Invalid field: {key}")
            setattr(self, key, value)
    @classmethod
    def from_row(cls, row):
        if len(row) != len(cls.fields):
            raise ValueError("Wrong number of fields")
        kwargs = dict(pair for pair in zip(cls.fields, row))
        return cls(**kwargs)
\end{codebox}

现在我可以为货运 CSV 文件格式创建一个具体的子类：

\begin{codebox}[python]
class DeliveryMapper(RowMapper):
    fields = ("destination", "method", "weight")
obj2 = DeliveryMapper.from_row(row1)
assert obj2.destination == "Sydney"
assert obj2.method == "truck"
assert obj2.weight == "25"
\end{codebox}

如果我要支持另一种 CSV 格式——比如搬家货车的物流——我只需提供列名，就能快速创建另一个子类：

\begin{codebox}[python]
class MovingMapper(RowMapper):
    fields = ("source", "destination", "square_feet")
\end{codebox}

虽然这样可行，但它不够 Pythonic。属性是用字符串而非变量名来指定的，这既让代码难以阅读，也让工具无所适从（见 Item 124（「考虑通过 \texttt{typing} 进行静态分析以消除缺陷」）和 Item 3（「永远不要指望 Python 在编译期检测错误」））。更重要的是，\texttt{fields} 元组与类体显得冗余：它是一份属性列表，却嵌套在另一份属性列表之中。

更好的做法是，我能把 CSV 列名直接写在类体里，像这样：

\begin{codebox}[python]
class BetterMovingMapper:
    source = ...
    destination = ...
    square_feet = ...
\end{codebox}

事实证明，这可以结合 Python 的三项特性来实现（另一种做法见 Item 51（「优先使用 \texttt{dataclasses} 定义轻量级类」））。第一项特性是 \texttt{\_\_init\_subclass\_\_} 这个特殊的类方法，它让你在定义一个子类时运行代码（见 Item 62（「用 \texttt{\_\_init\_subclass\_\_} 校验子类」））。第二项特性是：Python 的类属性可以在运行时通过类对象的 \texttt{\_\_dict\_\_} 实例字典来检查（见 Item 54（「考虑用 Mix-in 类组合功能」））。第三项特性是：Python 的字典会保留键/值对的插入顺序（见 Item 25（「依赖字典插入顺序时要谨慎」））。

这里我创建一个类，它会找出被赋值为 \texttt{...} 的子类属性，并把它们的名字存入 \texttt{fields} 类属性，供父类 \texttt{RowMapper} 使用：

\begin{codebox}[python]
class BetterRowMapper(RowMapper):
    def __init_subclass__(cls):
        fields = []
        for key, value in cls.__dict__.items():
            if value is Ellipsis:
                fields.append(key)
        cls.fields = tuple(fields)
\end{codebox}

现在我可以像之前一样声明一个具体类，但改用带省略号的类体来标明 CSV 文件的列：

\begin{codebox}[python]
class BetterDeliveryMapper(BetterRowMapper):
    destination = ...
    method = ...
    weight = ...
obj3 = BetterDeliveryMapper.from_row(row1)
assert obj3.destination == "Sydney"
assert obj3.method == "truck"
assert obj3.weight == "25"
\end{codebox}

如果 CSV 文件中列的顺序发生变化，我只需调整属性的定义顺序来补偿。例如，这里我把 \texttt{destination} 字段移到了末尾：

\begin{codebox}[python]
class ReorderedDeliveryMapper(BetterRowMapper):
    method = ...
    weight = ...
    destination = ...  # 已移动
row4 = ["road train", "90", "Perth"]  # 顺序不同
obj4 = ReorderedDeliveryMapper.from_row(row4)
print(obj4.__dict__)
>>>
{'method': 'road train', 'weight': '90', 'destination':
''Perth'}
\end{codebox}

在真实程序中，声明字段时我会使用描述符（descriptor）类而不是省略号，以支持诸如属性校验和数据转换之类的用例（背景知识见 Item 60（「用描述符实现可复用的 \texttt{@property} 方法」））。例如，假设我想让 \texttt{weight} 列被解析成浮点数，而不是保留为字符串。

这里我实现一个描述符（descriptor）类，它拦截属性访问，并按需转换被赋入的值：

\begin{codebox}[python]
class Field:
    def __init__(self):
        self.internal_name = None
    def __set_name__(self, owner, column_name):
        self.internal_name = "_" + column_name
    def __get__(self, instance, instance_type):
        if instance is None:
            return self
        return getattr(instance, self.internal_name, "")
    def __set__(self, instance, value):
        adjusted_value = self.convert(value)
        setattr(instance, self.internal_name, adjusted_value)
    def convert(self, value):
        raise NotImplementedError
\end{codebox}

我可以实现两个具体的 \texttt{Field} 子类——一个用于字符串，另一个用于浮点数：

\begin{codebox}[python]
class StringField(Field):
    def convert(self, value):
        if not isinstance(value, str):
            raise ValueError
        return value
class FloatField(Field):
    def convert(self, value):
        return float(value)
\end{codebox}

另一个用来表示 CSV 文件的新基类，可以改为查找 \texttt{Field} 实例（而非 \texttt{Ellipsis} 实例）来发现有序的 CSV 列，并相应地填充 \texttt{fields} 类属性：

\begin{codebox}[python]
class DescriptorRowMapper(RowMapper):
    def __init_subclass__(cls):
        fields = []
        for key, value in cls.__dict__.items():
            if isinstance(value, Field):  # 已修改
                fields.append(key)
        cls.fields = tuple(fields)
\end{codebox}

现在我可以为自己的特定 CSV 格式声明一个具体子类，而且 \texttt{weight} 字段会如预期那样被转换成浮点数：

\begin{codebox}[python]
class ConvertingDeliveryMapper(DescriptorRowMapper):
    destination = StringField()
    method = StringField()
    weight = FloatField()
obj5 = ConvertingDeliveryMapper.from_row(row1)
assert obj5.destination == "Sydney"
assert obj5.method == "truck"
assert obj5.weight == 25.0  # 数字，而非字符串
\end{codebox}

检查类属性也可以用来发现方法。与上面 CSV 用例完全不同，设想我想创建一个类来描述一组需要按定义顺序依次运行的方法所构成的有序工作流，像这样：

\begin{codebox}[python]
class HypotheticalWorkflow:
    def start_engine(self):
        ...
    def release_brake(self):
        ...
    def run(self):
        # 先运行 `start_engine` 再运行 `release_brake`
        ...
\end{codebox}

要使其可行，我首先创建一个简单的函数装饰器（function decorator）（见 Item 38（「用 \texttt{functools.wraps} 定义函数装饰器」）），用来标记哪些方法应当被纳入工作流：

\begin{codebox}[python]
def step(func):
    func._is_step = True
    return func
\end{codebox}

然后，一个新基类可以查找带有 \texttt{\_is\_step} 属性的可调用（callable）类属性（背景见 Item 48（「为简单接口接受函数而非类」）），从而发现哪些方法应当被纳入工作流，以及它们应当被调用的顺序：

\begin{codebox}[python]
class Workflow:
    def __init_subclass__(cls):
        steps = []
        for key, value in cls.__dict__.items():
            if callable(value) and hasattr(value, "_is_step"):
                steps.append(key)
        cls.steps = tuple(steps)
\end{codebox}

\texttt{run} 方法只需遍历步骤列表，并按保存的顺序调用这些方法。不需要其他任何样板代码：

\begin{codebox}[python]
    def run(self):
        for step_name in type(self).steps:
            func = getattr(self, step_name)
            func()
\end{codebox}

综合起来，这里我为启动汽车定义一个简单的工作流，其中包含了一个应当被基类忽略的辅助方法：

\begin{codebox}[python]
class MyWorkflow(Workflow):
    @step
    def start_engine(self):
        print("Engine is on!")
        ...
    def my_helper_function(self):
        raise RuntimeError("Should not be called")
    @step
    def release_brake(self):
        print("Brake is off!")
        ...
    ...
\end{codebox}

该工作流成功运行，并且不会调用那个有问题的（bad）方法：

\begin{codebox}[python]
workflow = MyWorkflow()
workflow.run()
>>>
Engine is on!
Brake is off!
...
\end{codebox}

\subsection{需要记住的要点}

\begin{itemize}
\item 你可以通过检查类对象对应的 \texttt{\_\_dict\_\_} 实例字典，在运行时查看类体中定义的属性和方法。
\item 类体的定义顺序会保留在类对象的 \texttt{\_\_dict\_\_} 中，使代码能够考虑类的属性和方法的相对位置。这在诸如把对象字段映射到 CSV 列索引这样的用例中尤其有用。
\item 描述符（descriptor）和方法装饰器（method decorator）可以用来进一步增强「利用类体定义顺序来控制程序（program）行为」这一做法的威力。
\end{itemize}

\section{Item 66: 组合式类扩展优先用类装饰器而非元类}

尽管元类（metaclass）能让你以多种方式定制类的创建过程（参见 Item 62（「使用 \texttt{\_\_init\_subclass\_\_} 校验子类」）和 Item 63（「使用 \texttt{\_\_init\_subclass\_\_} 登记类的存在」）），但它们仍然无法应对所有可能出现的情况。

举例来说，假设我想用一个辅助函数来装饰类的所有方法，该函数会打印参数、返回值以及抛出的任何异常。下面我定义这样一个调试装饰器（背景知识参见 Item 38（「使用 \texttt{functools.wraps} 定义函数装饰器」））：

\begin{codebox}[python]
from functools import wraps
def trace_func(func):
    if hasattr(func, "tracing"):  # 只装饰一次
        return func
    @wraps(func)
    def wrapper(*args, **kwargs):
        args_repr = repr(args)
        kwargs_repr = repr(kwargs)
        result = None
        try:
            result = func(*args, **kwargs)
            return result
        except Exception as e:
            result = e
            raise
        finally:
            print(
                f"{func.__name__}"
                f"({args_repr}, {kwargs_repr}) -> "
                f"{result!r}"
            )
    wrapper.tracing = True
    return wrapper
\end{codebox}

我可以把这个装饰器应用到新 \texttt{dict} 子类的各个特殊方法上（参见 Item 57（「从 \texttt{collections.abc} 类继承以实现自定义容器类型」））：

\begin{codebox}[python]
class TraceDict(dict):
    @trace_func
    def __init__(self, *args, **kwargs):
        return super().__init__(*args, **kwargs)
    @trace_func
    def __setitem__(self, *args, **kwargs):
        return super().__setitem__(*args, **kwargs)
    @trace_func
    def __getitem__(self, *args, **kwargs):
        return super().__getitem__(*args, **kwargs)
    ...
\end{codebox}

我可以通过操作该类的一个实例，来验证这些方法确实被装饰了：

\begin{codebox}[python]
trace_dict = TraceDict([("hi", 1)])
trace_dict["there"] = 2
trace_dict["hi"]
try:
    trace_dict["does not exist"]
except KeyError:
    pass  # 预期如此
>>>
__init__(({}, [('hi', 1)]), {}) -> None
__setitem__(({'hi': 1}, 'there', 2), {}) -> None
__getitem__(({'hi': 1, 'there': 2}, 'hi'), {}) -> 1
__getitem__(({'hi': 1, 'there': 2}, 'does not exist'),
'{}) -> KeyError('does not exist')
\end{codebox}

这段代码的问题在于，我必须把所有想用 \texttt{@trace\_func} 装饰的方法都重新定义一遍。这是冗余的样板代码，既难读又容易出错。此外，如果 \texttt{dict} 父类之后新增了某个方法，除非我也在 \texttt{TraceDict} 中定义它，否则它不会被装饰。

解决这个问题的一种方法是使用元类，自动装饰类的所有方法。下面我通过把新类型中的每个函数或方法都用 \texttt{trace\_func} 装饰器包装起来，来实现这一行为：

\begin{codebox}[python]
import types
TRACE_TYPES = (
    types.MethodType,
    types.FunctionType,
    types.BuiltinFunctionType,
    types.BuiltinMethodType,
    types.MethodDescriptorType,
    types.ClassMethodDescriptorType,
    types.WrapperDescriptorType,
)
IGNORE_METHODS = (
    "__repr__",
    "__str__",
)
class TraceMeta(type):
    def __new__(meta, name, bases, class_dict):
        klass = super().__new__(meta, name, bases, class_dict)
        for key in dir(klass):
            if key in IGNORE_METHODS:
                continue
            value = getattr(klass, key)
            if not isinstance(value, TRACE_TYPES):
                continue
            wrapped = trace_func(value)
            setattr(klass, key, wrapped)
        return klass
\end{codebox}

现在，我可以通过使用 \texttt{TraceMeta} 元类来声明我的 \texttt{dict} 子类，并验证它如预期般工作：

\begin{codebox}[python]
class TraceDict(dict, metaclass=TraceMeta):
    pass
trace_dict = TraceDict([("hi", 1)])
trace_dict["there"] = 2
trace_dict["hi"]
try:
    trace_dict["does not exist"]
except KeyError:
    pass  # 预期如此
>>>
__new__((<class '__main__.TraceDict'>, [('hi', 1)]), {}) -> {}
__init__(({}, [('hi', 1)]), {}) -> None
__setitem__(({'hi': 1}, 'there', 2), {}) -> None
__getitem__(({'hi': 1, 'there': 2}, 'hi'), {}) -> 1
__getitem__(({'hi': 1, 'there': 2}, 'does not exist'),
'{}) -> KeyError('does not exist')
\end{codebox}

这确实有效，它甚至打印出了我在先前实现中遗漏的对 \texttt{\_\_new\_\_} 的调用。如果某个父类已经指定了元类，而我仍尝试使用 \texttt{TraceMeta}，会发生什么呢？

\begin{codebox}[python]
class OtherMeta(type):
    pass
class SimpleDict(dict, metaclass=OtherMeta):
    pass
class ChildTraceDict(SimpleDict, metaclass=TraceMeta):
    pass
>>>
Traceback ...
TypeError: metaclass conflict: the metaclass of a derived class
'must be a (non-strict) subclass of the metaclasses of all its
'bases
\end{codebox}

这样会失败，是因为 \texttt{TraceMeta} 没有从 \texttt{OtherMeta} 继承。理论上，我可以通过让 \texttt{OtherMeta} 继承自 \texttt{TraceMeta} 来用元类继承解决此问题：

\begin{codebox}[python]
class TraceMeta(type):
    ...
class OtherMeta(TraceMeta):
    pass
class SimpleDict(dict, metaclass=OtherMeta):
    pass
class ChildTraceDict(SimpleDict, metaclass=TraceMeta):
    pass
trace_dict = ChildTraceDict([("hi", 1)])
trace_dict["there"] = 2
trace_dict["hi"]
try:
    trace_dict["does not exist"]
except KeyError:
    pass  # 预期如此
>>>
__init_subclass__((), {}) -> None
__new__((<class '__main__.ChildTraceDict'>, [('hi', 1)]),
'{}) -> {}
__init__(({}, [('hi', 1)]), {}) -> None
__setitem__(({'hi': 1}, 'there', 2), {}) -> None
__getitem__(({'hi': 1, 'there': 2}, 'hi'), {}) -> 1
__getitem__(({'hi': 1, 'there': 2}, 'does not exist'),
'{}) -> KeyError('does not exist')
\end{codebox}

但如果元类来自一个我无法修改的库，或者我想同时使用多个像 \texttt{TraceMeta} 这样的工具型元类，这种做法就行不通了。元类方案给被修改的类施加了过多的限制。

为了解决这个问题，Python 提供了类装饰器（class decorator）的支持。类装饰器的工作方式与函数装饰器完全一样：在类声明之前用一个函数加上 \texttt{@} 符号前缀来应用。该函数应当相应地修改或重新创建这个类，然后将其返回，如下所示：

\begin{codebox}[python]
def my_class_decorator(klass):
    klass.extra_param = "hello"
    return klass
@my_class_decorator
class MyClass:
    pass
print(MyClass)
print(MyClass.extra_param)
>>>
<class '__main__.MyClass'>
hello
\end{codebox}

我可以实现类装饰器，把上面 \texttt{TraceMeta.\_\_new\_\_} 方法的核心逻辑搬到一个独立函数中，从而把 \texttt{trace\_func} 函数装饰器应用到类的所有方法上。这个实现比元类版本短得多：

\begin{codebox}[python]
def trace(klass):
    for key in dir(klass):
        if key in IGNORE_METHODS:
            continue
        value = getattr(klass, key)
        if not isinstance(value, TRACE_TYPES):
            continue
        wrapped = trace_func(value)
        setattr(klass, key, wrapped)
    return klass
\end{codebox}

我可以把这个装饰器应用到我的 \texttt{dict} 子类上，从而得到与上面元类方案相同的行为：

\begin{codebox}[python]
@trace
class DecoratedTraceDict(dict):
    pass
trace_dict = DecoratedTraceDict([("hi", 1)])
trace_dict["there"] = 2
trace_dict["hi"]
try:
    trace_dict["does not exist"]
except KeyError:
    pass  # 预期如此
>>>
__new__((<class '__main__.DecoratedTraceDict'>, [('hi', 1)]),
'{}) -> {}
__init__(({}, [('hi', 1)]), {}) -> None
__setitem__(({'hi': 1}, 'there', 2), {}) -> None
__getitem__(({'hi': 1, 'there': 2}, 'hi'), {}) -> 1
__getitem__(({'hi': 1, 'there': 2}, 'does not exist'),
'{}) -> KeyError('does not exist')
\end{codebox}

当被装饰的类本身已经拥有元类时，类装饰器同样适用：

\begin{codebox}[python]
class OtherMeta(type):
    pass
@trace
class HasMetaTraceDict(dict, metaclass=OtherMeta):
    pass
trace_dict = HasMetaTraceDict([("hi", 1)])
trace_dict["there"] = 2
trace_dict["hi"]
try:
    trace_dict["does not exist"]
except KeyError:
    pass  # 预期如此
>>>
__new__((<class '__main__.HasMetaTraceDict'>, [('hi', 1)]),
'{}) -> {}
__init__(({}, [('hi', 1)]), {}) -> None
__setitem__(({'hi': 1}, 'there', 2), {}) -> None
__getitem__(({'hi': 1, 'there': 2}, 'hi'), {}) -> 1
__getitem__(({'hi': 1, 'there': 2}, 'does not exist'),
'{}) -> KeyError('does not exist')
\end{codebox}

当你在寻找可组合的方式来扩展类时，类装饰器就是完成这项任务的最佳工具。（参见 Item 104（「了解如何用 \texttt{heapq} 实现优先队列」），其中给出了一个名为 \texttt{functools.total\_ordering} 的类装饰器示例。）

\subsection{需要记住的要点}

\begin{itemize}
\item 类装饰器是一个简单的函数，它接收一个类实例作为参数，并返回一个新类或原始类的修改版本。
\item 当你想用最少的样板代码修改类的每个方法或属性时，类装饰器非常有用。
\item 元类不容易组合在一起，而多个类装饰器却可以无冲突地用于扩展同一个类。
\end{itemize}
